\chapter{LED e debug}
\label{ch:led}

\section{LED dell'Engineering Model}
Nessun LED è collegato a un GPIO dell'MCU \fonte{SCH}. I due LED danno un riscontro visivo indipendente dal firmware.

\begin{xltabular}{\linewidth}{@{}L{14mm}L{30mm}L{38mm}Y@{}}
\tabhead\hd{LED}&\hd{Pilotato da}&\hd{Acceso quando}&\hd{Note}\\
D203 (rosso) & \(\overline{\mathrm{RESET}}\) di IC205, tramite R209 \SI{1.5}{\kilo\ohm} & il Si4464 è fuori reset (SDN basso, modo A) & \SI{0.93}{\milli\ampere} continui. Spento durante il reset della radio o con il 3V3 sotto soglia. Nel modo B non indica più lo stato della radio \fonte{LOG rev.\ 1.45, P14; §46}.\\
D9 (verde) & 3V3 tramite R49 \SI{10}{\kilo\ohm} e JP1 & 3V3 presente \textbf{e} JP1 chiuso & JP1 aperto di serie. Luce molto debole (circa \SI{0.09}{\milli\ampere}), voluto: serve solo per il debug iniziale \fonte{LOG S-13}.\\
\bottomrule
\end{xltabular}

\begin{nota}[come usare D203 al banco]
D203 è il modo più rapido per vedere un reset della radio durante la trasmissione (calo del 3V3, misura M-03): se si spegne per circa \SI{200}{\milli\second} durante una TX, IC205 è intervenuto. In parallelo il firmware deve rilevare l'evento con F-10.
\end{nota}

\section{Flight Model}
\begin{regola}[LED sul Flight Model]
Sul Flight Model i LED di debug non si montano: D203 e R209 (e D9) vanno esclusi dalla BOM di volo \fonte{LOG rev.\ 1.45, P14}. Il firmware non deve dipendere dai LED. Ogni informazione utile che oggi si legge dai LED va anche in telemetria (stato della radio, reset, eventi F-10).
\end{regola}

Altre differenze tra configurazione di sviluppo e di volo:
\begin{xltabular}{\linewidth}{@{}L{30mm}L{46mm}Y@{}}
\tabhead\hd{Elemento}&\hd{Sviluppo (banco)}&\hd{Volo}\\
R413 (EN watchdog) & montata: watchdog spento & \textbf{non montata}: watchdog attivo\\
R205 / R204 & modo A di serie; modo B solo se M-03 fallisce & secondo l'esito di M-03 (costante \code{MODO\_A/B})\\
LED D203, D9 & montati & non montati\\
JP201, JP6 & secondo il banco & secondo l'architettura del bus del satellite\\
Firmware & build di collaudo con comandi di test & build di volo: PA\_PWR\_LVL fisso (F-01), nessun comando di test RF\\
\bottomrule
\end{xltabular}

\section{Strumenti di debug}
\begin{itemize}
\item \textbf{SWD} su J1 (TC2030): SWDIO, SWCLK, NRST, SWO, 3V3 e GND. Serve R413 montata per lavorare con i breakpoint.
\item \textbf{Console} su USART1 tramite CH340E e USB-C. La scheda deve avere il 3V3 (EN\_DOCK).
\item \textbf{Bootloader di sistema}: TP1 (BOOT0) alto al reset, STM32CubeProgrammer in UART \fonte{LOG F-07}.
\item \textbf{Test point}: TP701 (XIN Si4464, clock del TCXO), TP702 (PC2), TP703 (EN\_RX\_LNA), TP1 (BOOT0) \fonte{SCH}.
\end{itemize}
